% Preliminary arithmetic proof note; the complete integrated proof is now
% sections/C_gz_factors.tex, Statement C.6.R5. Not included automatically.
\section*{A denominator certificate for the six remaining C.6 factors}

The primary source is Cai, Shu and Tian, \emph{Explicit Gross--Zagier and
Waldspurger formulae}, Algebra \& Number Theory 8(10) (2014), 2523--2572,
DOI 10.2140/ant.2014.8.2523. The complete publisher PDF was opened at
\url{https://msp.org/ant/2014/8-10/ant-v8-n10-p05-p.pdf}.
The exact hypotheses and formula used below are Theorem 1.1 on
pp.~2524--2525; the diagonal height is the N\'eron--Tate pairing over
$K$, as in the general formulation on p.~2531.

\paragraph{Quadratic twisting of a known parametrization.}
Let $E_0/\Q$ have conductor $N_0$, and let
$\phi:X_0(N_0)\to E_0$ send $\infty$ to $O$ and satisfy
$\phi^*\omega_0=c_0 f_0(q)dq/q$. Let $b$ be an odd fundamental
discriminant coprime to $N_0$, put $m=|b|$, and assume that the minimal
twist $E_0^b$ has conductor $N_0b^2$ and
$c_4(E_0^b)=b^2c_4(E_0)$. The last equality says that the twist
requires no scaling of the minimal differential. Then there is a
parametrization $X_0(N_0b^2)\to E_0^b$ with Manin factor $c_0$.

Indeed, using the group law on $E_0$, put
\[
 F(z)=\sum_{a\bmod m}\chi_b(a)\phi(z+a/m).
\]
For $\gamma=\left(\begin{smallmatrix}A&B\\C&D\end{smallmatrix}\right)
\in\Gamma_0(N_0m^2)$, choose $a'\equiv aD^2\pmod m$. Direct
matrix multiplication shows
$T_{a/m}\gamma T_{a'/m}^{-1}\in\Gamma_0(N_0)$.
Because $\chi_b(D^2)=1$, this proves $F(\gamma z)=F(z)$.
The summands are algebraic maps on a common congruence cover, so their
sum descends to an algebraic map on $X_0(N_0m^2)$ over
$\Q(\zeta_m)$. Its expansion at the rational cusp $\infty$ shows
$F^{\sigma_u}=\chi_b(u)F$, where
$\sigma_u(\zeta_m)=\zeta_m^u$; hence quadratic descent gives a map
to $E_0^b$ over $\Q$. It sends $\infty$ to $O$. Finally
\[
 F^*\omega_0=c_0\sum_a\chi_b(a)f_0(z+a/m)\,2\pi i\,dz
 =c_0\tau(\chi_b)f_b(q)\frac{dq}{q}.
\]
Here $\tau(\chi_b)^2=b$. Under the twist isomorphism
$\iota:E_0^b\to E_0$, the minimal differentials satisfy
$\iota^*\omega_0=\sqrt b\,\omega_b$. Choosing the square root
consistent with the Gauss sum proves the assertion. A sign has no
effect on the positive Manin factor. This argument proves what is
needed for the particular twists, without using a general Manin
constant conjecture.

If $c_0=1$ and $t_b=|E_0^b(\Q)_{\rm tors}|$, it follows that
\[
 t_b n_b\,\frac{L(E_0^b,1)}{\Omega_b}\in\Z,
 \qquad n_b=\#\pi_0(E_0^b(\R))\in\{1,2\}.
\]
To see this, the cuspidal divisor $[0]-[\infty]$ maps to a rational
torsion point, by Manin--Drinfeld. Its real elliptic logarithm is
$\pm L(E_0^b,1)$. Multiplication by $t_b$ gives a real period, and
$\Omega_b=n_b\Omega_b^+$.

\paragraph{Genus characters at the small base level.}
For each row use the following data. The base Manin factors are all one. For the first five bases this was
independently returned by the unconditional PARI modular-symbol algorithm;
for the last base it is also explicitly recorded in Belabas, Bernardi and
Perrin-Riou (2019), the unnumbered 26569a1/26569a2 block in
Appendix~B, p.~102, for 26569a1. That source was opened in full
at \url{https://www.numdam.org/item/10.5802/pmb.37.pdf}; its Algorithm~1
explains the computation. All equations, conductors,
twist identities, differential scalings, torsion orders and splitting
conditions are checked exactly in \path{code/B_C6_gz_exact_metadata.gp}.
\[
\begin{array}{c|c|r|r|r}
 i&E_0&N_0&d&D=\operatorname{disc}(K)\\\hline
 2&576\mathrm e2&576&-11&-23\\
 4&3136\mathrm r2&3136&57&-31\\
 6&121\mathrm b1&121&21&-19\\
 9&1849\mathrm a1&1849&21&-19\\
 10&4489\mathrm a1&4489&217&-3\\
 11&26569\mathrm a1&26569&185801&-3
\end{array}
\]
In every case $E_i=E_0^d$, $d,D,Dd$ are fundamental discriminants,
$\gcd(dD,N_0)=\gcd(d,D)=1$, and every prime dividing $N_0$ splits
in $K=\Q(\sqrt D)$. The genus character of
$K(\sqrt d)/K$ is a primitive character of
$\operatorname{Pic}(\mathcal O_{|d|})$, of conductor $|d|$.
For example it is given on ideals prime to $d$ by
$\chi(\mathfrak a)=\chi_d(N\mathfrak a)$; it is trivial on the
units of the order because their norms are squares modulo $d$.
At every prime dividing $d$ the local norm character is nontrivial,
so the conductor is exactly $|d|$. All Heegner hypotheses of CST
Theorem 1.1 hold. Its $\mu(N_0,D)$ is zero, and its
$u=[\mathcal O_{|d|}^\times:\Z^\times]$ is one, including
$D=-3$, since $|d|>1$.

Let $Y=\sum_{[\mathfrak a]}\chi(\mathfrak a)\phi(P_{\mathfrak a})$.
This is an actual point, since the coefficients are $\pm1$.
It is fixed by $\ker\chi$ and is anti-invariant for the other
automorphism of $K(\sqrt d)/K$, so twisting gives
$Y_d\in E_i(K)$. Artin factorization gives
\[
 L(s,E_0,\chi)=L(E_i,s)L(E_i^D,s).
\]
The auxiliary central values below are nonzero. Consequently
$E_i^D(\Q)$ has rank zero, and $E_i(K)$ has rank one. Since the
specified $P_i$ generates $E_i(\Q)$ modulo torsion, trace gives
\[
 2Y_d=aP_i\pmod{E_i(K)_{\rm tors}},\qquad a\in\Z.
\]
Indeed, $Y_d-\overline{Y_d}$ is torsion, while
$Y_d+\overline{Y_d}$ is rational. Therefore
$h_O(Y_d)=a^2h_O(P_i)/4$.

\paragraph{Normalization of the Gross--Zagier formula.}
Write $V_0=\operatorname{area}(\C/\Lambda_0)$ for the area determined
by the minimal differential. Integration of $\phi^*\omega_0$ gives
\[
 4\pi^2(f_0,f_0)_{\Gamma_0(N_0)}
 =\frac{\deg\phi}{c_0^2}V_0.
\]
Theorem 1.1 consequently gives
\[
 L'(E_i,1)L(E_i^D,1)
 =\frac{2V_0}{c_0^2|d|\sqrt{|D|}}\,\widehat h_K(Y).
\]
For clarity, this normalization follows from the Poincar\'e bundle,
not from the choice of a numerical height routine. Let $L=\mathcal O(O)$
and let $\lambda_L:E\to E^\vee$ be its principal polarization.
The pullback of the normalized Poincar\'e bundle to $E\times E$
is $m^*L\otimes p_1^*L^{-1}\otimes p_2^*L^{-1}$. Its restriction
to the diagonal is $[2]^*L\otimes L^{-2}\simeq L^2$ for canonical
heights. Hence the diagonal Poincar\'e height is $2h_O$. Relative
heights over $K$ multiply the absolute height by $[K:\Q]$;
therefore CST's diagonal N\'eron--Tate pairing is
$2[K:\Q]h_O=4h_O$. Thus the last height is
$4h_O(Y_d)=a^2h_O(P_i)$. The product of the two full real periods is
\[
 \Omega_i\Omega_{i,D}=
 \begin{cases}
 2V_0/(|d|\sqrt{|D|}),&\Delta(E_0)<0,\\
 4V_0/(|d|\sqrt{|D|}),&\Delta(E_0)>0.
 \end{cases}
\]
This follows directly from a real-adapted lattice basis; both
minimal twist differential scalings are one. Put
$r_i=L(E_i^D,1)/\Omega_{i,D}$ and
$q_i=L'(E_i,1)/(\Omega_i h_O(P_i))$. Then
\[
 q_ir_i=\begin{cases}
 a^2/c_0^2,&\Delta(E_0)<0,\\
 a^2/(2c_0^2),&\Delta(E_0)>0.
 \end{cases}
\]
In particular, if $c_0=1$ and $r_i=A_i/B_i$ is in lowest terms,
the denominator of $q_i$ divides $2|A_i|$. For the numerical
separation certificate we may safely use the weaker bound
$4|A_i|$.

\paragraph{Completed computational certificate.}
The auxiliary ratios are $64,64,8,16,400,64$ in the displayed row
order. The exact $q_i$ are $8,32,8,16,16,576$. The corresponding
absolute integers $|a|$ are $32,64,8,16,80,192$, respectively.
The complete certificate is \path{code/B_C6_exact_factors.py},
with all input and output intervals in
\path{records/B_C6_exact_factors.json}. Both $L$-sums use MPFR
directed rounding, explicit Taylor remainder bounds and infinite
tails; periods use Arb with an explicit Eisenstein-series tail;
heights use the independent \path{records/C_cm_height_arb.json}.
The primary YZZ manuscript, Section 7.1.1, printed pp.~218--219,
Theorem 7.1.2, has now been read and explicitly confirms the
relative-field and Poincar\'e conventions used above. The full
computational proof, including global identification of the CM
coefficient streams through the valence bound, is in
\path{sections/C_gz_factors.tex}, Statement C.6.R5.
